\section{Introduction}
\label{sec:intro}

Large language models (LLMs) increasingly run close to their users, on phones, robots and embedded boards, for responsiveness, privacy and cost. Such devices can host only small models, while the quality users expect comes from models with billions of parameters running in the cloud. Speculative decoding~\cite{leviathan2023,chen2023} offers a natural split between the two. A small draft model proposes several tokens, and the large target model checks all of them in one forward pass. With greedy verification every emitted token is the target model's own greedy choice, yet each target pass can emit several tokens. Placing the draft model on the device and the target model in the cloud combines local responsiveness with cloud-level quality, and several recent systems explore this split~\cite{sled2025,specedge2025,dsd2025,picospec2026,wisp2026}.

The split exposes a cost that single-machine speculative decoding does not have: every verification is a network round trip. In a synchronous edge--cloud loop, the device drafts, sends the draft and waits idle for the verdict before drafting again. A natural remedy is \emph{lookahead}. While the cloud verifies the current block, the device already drafts the next one, assuming that the current block will be accepted; parallel drafting and verification has been explored on a single server~\cite{pearl2025}. Lookahead turns idle waiting into useful work, but it creates two tensions that matter at the edge. First, the speculated continuation is usable only if the verifier accepts the entire block \emph{and} its bonus token matches the draft's guess. Otherwise it is discarded, and on an energy-constrained device every discarded token is wasted energy. Second, the right amount of lookahead depends on quantities that change across requests and rounds: the network delay, the verification time, the draft speed and how predictable the text currently is.

We study these tensions on a real deployment: an NVIDIA Jetson Orin Nano drafting with Llama-3.2-1B over campus Wi-Fi for a Llama-3.1-8B verifier served by vLLM~\cite{vllm2023} on an RTX 3080 server, with the board's input power measured for every request. Three observations shape our design (Section~\ref{sec:measure}):
\begin{itemize}
\item \textbf{Drafting sits on the critical path.} In synchronous speculative decoding, local drafting takes \E{sync:low:share} of a round with 2\,ms of added delay and \E{sync:high:share} with 100\,ms.
\item \textbf{Verification cost is almost flat in the draft length.} Each additional draft token adds only \E{verify:draft}\,ms to a mean verification time of \E{verify:mean}\,ms, because the cost is dominated by recomputing the context.
\item \textbf{Draft work has a predictable energy price.} Per-request board energy is explained by wall time and draft-server busy time ($R^2=\E{energy:rsq}$), so wasted drafting translates directly into joules.
\end{itemize}
Together, these observations suggest a simple principle. \emph{Tokens drafted after a verification returns delay the next request, whereas tokens drafted during verification are nearly free in time but not in energy.} A good controller should therefore draft as little as possible on the critical path, and draft ahead only as much as is likely to be reused.

We present \cape{} (Critical-path-Aware PrEfetching), an edge-only controller built on this principle. \cape{} models the decoding pipeline as a two-state Markov chain whose states record whether the current round can reuse a prefetched block (a hit) or must draft afresh (a miss). For every pair of fresh draft window $\gamma_f$ and lookahead length $h$, the model predicts the long-run time, draft-busy time and number of emitted tokens per round. \cape{} picks the pair that minimizes an energy-weighted cost per token, where a single weight $w$ trades speed against energy. Its inputs are estimated online from the last few rounds: state-dependent acceptance rates, the verification round trip and the per-token draft cost. \cape{} keeps strict greedy verification, so every emitted token is the target model's top choice, and it requires no change to the cloud server. The controller adds 170 lines of Python to a lookahead client and decides in about \E{overhead:ms}\,ms per round.

We evaluate \cape{} on all \E{full:prompts} prompts of our GSM8K, HumanEval and MT-Bench sets at two network delays, against synchronous decoding, three fixed windows and AdaEDL~\cite{adaedl2024}: \E{full:requests:fmt} measured requests with randomized interleaving and paired prompt-cluster bootstrap confidence intervals. With its speed setting, \cape{} is \E{full:capes:fixed4:pooled:abssp} \E{full:capes:fixed4:pooled:spci} faster than the fastest fixed window, and faster on every task. It even beats an oracle that picks the best fixed window for each delay, by \E{full:capes:fixed4:low:abssp} at 2\,ms and \E{full:capes:fixed8:high:abssp} at 100\,ms, while using less energy per token than that oracle at both delays. With its energy setting, \cape{} keeps the speed of the fastest fixed window and lowers J/token by \E{full:capee:fixed4:pooled:abssv} \E{full:capee:fixed4:pooled:svci}. It is \E{full:capee:sync4:pooled:abssp}--\E{full:capes:sync4:pooled:abssp} faster than synchronous speculative decoding and \E{full:capee:adaedl:pooled:abssp}--\E{full:capes:adaedl:pooled:abssp} faster than AdaEDL, and its energy setting is both faster and more energy-efficient than AdaEDL and an eight-token fixed window. On two earlier held-out prompt sets, the default setting was \E{h1:cape:fixed4:pooled:abssp} and \E{h2:cape:fixed4:pooled:abssp} faster than the fixed window without increasing energy per token.

In summary, this paper makes the following contributions:
\begin{itemize}
\item A measurement study of edge--cloud speculative decoding on an embedded GPU. It quantifies the critical-path share of drafting and the near-constant cost of verification, and it fits a per-request board energy model (Section~\ref{sec:measure}).
\item \cape{}, a critical-path-aware lookahead controller based on a two-state Markov model with an explicit energy price. It runs entirely on the device and keeps strict greedy verification (Section~\ref{sec:design}).
\item An evaluation on real hardware with board-level energy measurement: \E{full:requests:fmt} requests on \E{full:prompts} prompts against synchronous decoding, three fixed windows and AdaEDL, plus two held-out prompt sets and a heuristic latency-regime controller (Section~\ref{sec:eval}).
\end{itemize}
